\documentclass{article}
\usepackage{graphicx} % Required for inserting images
\usepackage{xcolor}
\usepackage[utf8]{inputenc}
\usepackage[spanish]{babel} % Para que diga "Índice" en vez de "Contents"
\usepackage{hyperref} % Hace que el índice tenga links clicables
\hypersetup{
    colorlinks=true,
    linkcolor=black, % Cambia a 'blue' si quieres que los links se vean azules
    filecolor=magenta,      
    urlcolor=cyan,
}

\begin{figure}[]
    \centering
    \includegraphics[width=0.6\textwidth]{download.png}
    \label{fig:UMS_logo}
\end{figure}

\title{Taller Programación: Funnys Company}
\author{FREE \$BTC me hackeo ELON MUSK }
\date{Marzo 2026}

\begin{document}

\maketitle
\newpage

\tableofcontents
\newpage

\section{Introducción}
Funnys Company es una empresa orientada a la producción de artículos para la entretención en el hogar. Actualmente opera con una planta pequeña ubicada en Rancagua. Recientemente, el mercado experimentó fuertes cambios que resultaron en un aumento de la demanda nacional.

El problema principal es que la capacidad de producción actual y la red de distribución existente son insuficientes para satisfacer este nuevo escenario, lo que impide a la empresa satisfacer la demada. Por ello, la administración nos ha contactado para rediseñar estratégicamente su cadena de suministro para los proximos tres años.

Dentro de las alternativas de solución evaluadas, se contempla la posibilidad de expandir la capacidad mediante la apertura de nuevas plantas de producción (de tamaño pequeño o grande) seleccionando entre cinco ciudades candidatas, además de determinar la combinación óptima entre tres tipos de transporte (AT1, AT2 y AT3) para llevar los productos a seis regiones del país.

Para resolver el problema, nuestro equipo propone implementar un modelo de optimización matemática lineal entera mixta. Dado que no se cuenta con información sobre los precios de venta para maximizar ganancia, el enfoque será la minimización de los costos. Utilizando el solver Gurobi con Python, el modelo buscará la configuración exacta de la red (dónde y qué tipo de plantas abrir, qué rutas utilizar cada año) que garantice la satisfacción total de la demanda al menor costo posible, equilibrando los costos fijos con los costos variables de producción y transporte.

Nosotros proponemos resolver este problema de minimización con un mayor enfoque en decidir el tipo de transporte a utilizar entre cada planta y región según cada año de la planificación. 

\section{Modelo Matematico}
Para abordar el problema, se formuló un modelo de Programación Lineal Entera Mixta (MILP). Este modelo busca minimizar los costos totales de la cadena de suministro decidiendo la ubicación y tamaño de nuevas plantas (variables binarias) y la cantidad de productos a enviar por cada tipo de transporte (variables enteras). A continuación, se detallan los componentes del modelo:
\subsection{Conjuntos}
\begin{itemize}
    \item $C$: Ciudades de origen $\{\text{Antofagasta, Valparaíso, Santiago, Concepción, Puerto Montt, Rancagua}\}$
    \item $P$: Tipos de plantas de producción $\{\text{Pequeña, Grande}\}$
    \item $R$: Regiones de destino $\{1, 2, 3, 4, 5, 6\}$
    \item $T$: Tipos de transporte $\{\text{AT1, AT2, AT3}\}$
    \item $A$: Años $\{1, 2, 3\}$
\end{itemize}

\subsection{Parámetros}
\begin{itemize}
    \item $D_r$: Demanda base de la región $r$ en el año 1
    \item $Tasa_r$: Tasa de crecimiento anual de la demanda en la región $r$
    \item $Cap_p$: Capacidad máxima de producción de una planta tipo $p$
    \item $CF_{cp}$: Costo fijo anual de operación de una planta tipo $p$ en la ciudad $c$
    \item $CA_{cp}$: Costo de apertura de una planta tipo $p$ en la ciudad $c$
    \item $CV_c$: Costo variable de producción por unidad en la ciudad $c$ (independiente del tipo de planta)
    \item $TRA_{tcr}$: Costo unitario de transporte desde la ciudad $c$ a la región $r$ usando el medio $t$
\end{itemize}

\subsection{Variables de Decisión}
\begin{itemize}
    \item $X_{crta}$: Cantidad de unidades enviadas desde la ciudad $c$ a la región $r$, usando el transporte $t$, durante el año $a$.
    \item $Y_{cp}$: Toma el valor de 1 si se instala una planta tipo $p$ en la ciudad $c$, y 0 en caso contrario.
\end{itemize}

\subsection{Función Objetivo}
$$ \min Z = \sum_{c \in C} \sum_{p \in P} \left( CA_{cp} + 3 \cdot CF_{cp} \right) \cdot Y_{cp} + \sum_{c \in C} \sum_{r \in R} \sum_{t \in T} \sum_{a \in A} \left( CV_c + TRA_{tcr} \right) \cdot X_{crta} $$

\subsection{Restricciones}
\begin{itemize}
    \item \textbf{Condición inicial (Planta existente en Rancagua):} Ya existe una planta pequeña en Rancagua por ende es necesario que $Y_{cp}$ sea 1 ahí.
    $$ Y_{\text{Rancagua, Pequeña}} = 1 $$
    
    \item \textbf{Satisfacción de la demanda:} Se necesita que la cantidad transportada para cierta región sea mayor a la demanda de esa región durante todos los años.
    $$ \sum_{c \in C} \sum_{t \in T} X_{crta} \geq D_r \cdot (1 + Tasa_r)^{a-1} \quad \forall r \in R, \forall a \in A $$
    
    \item \textbf{Límite de instalación por ciudad:} En una ciudad no puede existir mas de una planta.
    $$ \sum_{p \in P} Y_{cp} \leq 1 \quad \forall c \in C $$

    \item \textbf{Límite de cantidad de plantas pequeñas a nivel nacional:} Funnys company solo puede tener 1 planta pequeña a nivel nacional, sin contar la planta en Rancagua, por ende limitamos la apertura de plantas pequeñas.
    $$ \sum_{c \in C} Y_{c, \text{Pequeña}} \leq 2 $$

    \item \textbf{Límite de cantidad de plantas grandes a nivel nacional:} Funnys company solo puede tener 1 planta grande a nivel nacional, por ende limitamos la apertura de plantas grandes.
    $$ \sum_{c \in C} Y_{c, \text{Grande}} \leq 1 $$
    
    \item \textbf{Balance de producción y capacidad:} No se puede mandar elementos desde una planta, si no esta creada.
    $$ \sum_{r \in R} \sum_{t \in T} X_{crta} \leq \sum_{p \in P} Cap_p \cdot Y_{cp} \quad \forall c \in C, \forall a \in A $$
    
    \item \textbf{Naturaleza de las variables:}
    $$ X_{crta} \in \mathbb{Z}^+ \cup \{0\} \quad \forall c \in C, \forall r \in R, \forall t \in T, \forall a \in A $$
    $$ Y_{cp} \in \{0, 1\} \quad \forall c \in C, \forall p \in P $$
\end{itemize}

\section{Supuestos}
Para la construcción del modelo, se han establecido los siguientes supuestos que permiten establecer limites iniciales:
\begin{itemize}
    \item \textbf{Instalación en el Año 1}: Cualquier planta nueva que el modelo decida abrir, se abre en el año 1 y opera durante los 3 años. El modelo no permite abrir una planta en el año 2 o en el año 3.
    \item \textbf{Apertura en Rancagua}: Como Rancagua ya existe, se obliga su existencia con una planta pequeña.
    \item \textbf{límite de plantas a nivel nacional}: Se pueden instalar máximo 2 plantas pequeñas y 1 grande en todo el país. Son 2 pequeñas por que se cuenta la de Rancagua como una de estas.
    \item \textbf{Divisibilidad del transporte}: El modelo asume que se puede enviar cualquier cantidad entera de productos (ej. 1, 15, o 900.000 unidades) y que el costo es estrictamente proporcional por unidad. Asume que no hay "camiones a medio llenar" que cobren una tarifa fija por viaje, sino un costo unitario puro.
    \item \textbf{Análisis transporte anual}: El modelo se encarga de optimizar el costo de los transportes de los productos durante cada año, por eso se ocupa el indice de año solamente en este.
    \item \textbf{Decisión transporte}: No se decide por solo un tipo de transporte que sea más conveniente en general, sino el mejor transporte para cada ciudad.
    \item \textbf{Costos variables}: Como los costos variables son iguales sin importar el tipo de planta, el parametro "CV" no contiene la planta.
    \item \textbf{Producción Rancagua}: Rancagua empieza a producir desde el año 1, no tiene producción anterior. Solo se ahorra la apertura de la planta.
    \item \textbf{Producción inmediata}: Se asume que la instalación de una planta será instantánea, por ende puede comenzar a producir inmediatamente.
    \item \textbf{Ubicación regiones}: Se hizo una traducción de R1, R2 y demás regiones a como se conocen generalmente en Chile, por ejemplo, R5 es la región de Valparaíso.
    \item \textbf{La bencina no sufre de inflación}: El modelo considera que los precios por transportar productos no cambian con el paso del tiempo
    
\end{itemize}

\section{Configuración optima}

Basado en los resultados obtenidos, la configuración óptima recomendada para Funnys Company consiste en una red de dos plantas de capacidad pequeña:

\begin{itemize}
\item \textbf{Planta Rancagua (Existente):} Mantiene su operación como la planta principal de distribución, operando a máxima capacidad a partir del segundo año.
\item \textbf{Planta Antofagasta (Nueva):} Se recomienda una planta pequeña en Antofagasta. Aunque la apertura se realiza en el Año 1, su utilización comienza en el Año 2 para absorber el crecimiento de la demanda en la zona norte.
\end{itemize}

A continuación, se detalla el flujo de productos y la selección de transporte para los tres proximos años:

\begin{table}[h!]
\centering
\begin{tabular}{lrrrrrrr}
\hline
Ciudad & R1 & R2 & R3 & R4 & R5 & R6 & Total \\
\hline
Antofagasta & 0 & 0 & 0 & 0 & 0 & 0 & 0 \\
Rancagua & 951776 & 967364 & 512051 & 386248 & 946174 & 303445 & 4067058 \\
\hline
Total & 951776 & 967364 & 512051 & 386248 & 946174 & 303445 & 4067058 \\
\hline
\end{tabular}
\caption{Transporte Año 1}
\end{table}

\begin{table}[h!]
\centering
\begin{tabular}{lrrrrrrr}
\hline
Ciudad & R1 & R2 & R3 & R4 & R5 & R6 & Total \\
\hline
Antofagasta & 446832 & 0 & 0 & 0 & 0 & 0 & 446832 \\
Rancagua & 657229 & 1180185 & 645185 & 444186 & 1315182 & 394479 & 4636446 \\
\hline
Total & 1104061 & 1180185 & 645185 & 444186 & 1315182 & 394479 & 5083278 \\
\hline
\end{tabular}
\caption{Transporte Año 2}
\end{table}

\begin{table}[h!]
\centering
\begin{tabular}{lrrrrrrr}
\hline
Ciudad & R1 & R2 & R3 & R4 & R5 & R6 & Total \\
\hline
Antofagasta & 1280710 & 468051 & 0 & 0 & 0 & 0 & 1748761 \\
Rancagua & 0 & 971774 & 812933 & 510813 & 1828103 & 512823 & 4636446 \\
\hline
Total & 1280710 & 1439825 & 812933 & 510813 & 1828103 & 512823 & 6385207 \\
\hline
\end{tabular}
\caption{Transporte Año 3}
\end{table}

\subsection{Estrategia de Transporte}
La selección del transporte se basa en la minimización del costo unitario por ruta. Dado que el modelo no presenta restricciones de flota, la elección es constante en el tiempo para cada par Ciudad-Región:
\begin{table}[h!]
\centering
\begin{tabular}{lcccccc}
\hline
Ciudad & R1 & R2 & R3 & R4 & R5 & R6 \\
\hline
Antofagasta & AT1 & AT1 & - & - & - & - \\
Rancagua    & AT2 & AT1 & AT2 & AT1 & AT3 & AT1 \\
\hline
\end{tabular}
\caption{Tipo de transporte por ruta (Ciudad $\rightarrow$ Región)}
\end{table}

\subsection{Análisis de utilización de red}
Es importante destacar que bajo la estructura actual del modelo, la planta de Antofagasta presenta nada de producción durante el Año 1. Esto ocurre porque la capacidad de la planta de Rancagua es suficiente para cubrir la demanda nacional inicial a un costo operativo menor. La apertura de Antofagasta se justifica para el Año 2 y 3, donde se vuelve necesaria para satisfacer el aumento de la demanda.

\subsection{Visualización}
Para complementar el análisis numérico, se presenta la configuración de la red logística de Funnys Company distribuida en los tres años. En la Figura \ref{fig:mapa_chile}, se pueden observar los siguientes elementos:

\textbf{Plantas:} Se representan con figuras cuadradas. El color \textbf{naranja} indica una planta \textbf{Pequeña}, mientras que el \textbf{rojo} es para plantas \textbf{Grandes}, aunque en esta ocasión no se logre ver una.

\textbf{Flujos del transporte:} Las flechas indican el movimiento de productos desde las plantas hacia las seis regiones. El \textbf{grosor de la línea} es directamente proporcional al volumen de carga transportada, permitiendo identificar visualmente las rutas críticas de la operación.

\textbf{Codificación por Transporte:} Los colores de las rutas diferencian el tipo de servicio seleccionado por el modelo para optimizar costos:
\begin{itemize}
\item \textcolor{blue}{\textbf{Azul}}: Uso del servicio de transporte \textbf{AT1}.
\item \textcolor{green}{\textbf{Verde}}: Uso del servicio de transporte \textbf{AT2}.
\item \textcolor{purple}{\textbf{Morado}}: Uso del servicio de transporte \textbf{AT3}.
\end{itemize}

Esta visualización permite confirmar cómo la planta de Antofagasta, aunque habilitada desde el inicio, comienza su despliegue logístico a partir del segundo año.
\begin{figure}[h!]
    \centering
    \includegraphics[width=0.7\textwidth]{mapa_logistica_funnycompany.png}
    \caption{Visualización geográfica de la red de transporte para Funnys Company.}
    \label{fig:mapa_chile}
\end{figure}

\section{Relajación Restricción}
Al relajar la restricción de unicidad de instalaciones por ciudad y permitir que la variable $Y_{cp}$ sea entera (pudiendo tomar valores $\geq 1$), el modelo experimenta un cambio significativo en su estructura de red.
\begin{itemize}
    \item \textbf{Solución:} La solución óptima se desplaza hacia la concentración de planta en una sola ciudad. El modelo decide instalar dos plantas pequeñas en Rancagua.
    \item \textbf{Justificación:} Este cambio se debe a que Rancagua presenta las condiciones de costo más favorables (inexistente costo de apertura, menores costos variables y de transporte combinados para la mayoría de las regiones).
\end{itemize}
\subsection{Análisis Adicional}Como ejercicio complementario, se evaluó un escenario donde se cambia el límite de plantas nacionales, 1 grande y 1 pequeña:
\begin{itemize}
    \item El modelo deja la pequeña en Rancagua.
    \item Abre una planta grande en Valparaíso.
\end{itemize}

Con estas condiciones Valparaiso toma el rol de cumplir la demanda de R1 y R2 que tiene Antofagasta en la solución original.

\section{Conclusiones}
Tras la ejecución del modelo y el análisis de los resultados, se extraen las siguientes conclusiones:

\begin{itemize}

    \item \textbf{Impacto estructural de la localización}: La selección de las ciudades es la variable binaria que determina la estructura de costos del modelo. Al fijar estas variables, se condicionan inmediatamente los demás parámetros de la función objetivo: costos fijos, variables y de transporte. Esto convierte a la localización en la restricción principal que define la calidad de la solución final.
    
    \item \textbf{Limitaciones de la solución}: La solución obtenida es un óptimo atada a los supuestos actuales del modelo. Al forzar la apertura de todas las plantas en el Año 1, la función objetivo asume una penalización monetaria por los costos fijos de Antofagasta, a pesar de no tener una producción. Para acercarse a un óptimo global más realista, sería necesario agregar un índice temporal a la variable de apertura ($Y_{cpa}$), permitiendo al algoritmo evaluar si diferir la apertura al año 2 genera una función objetivo de menor valor.

    \item \textbf{Complejidad dimensional por rutas de transporte}: Mientras la localización fija los costos, la asignación de transporte es la principal responsable de la complejidad computacional del problema. Al existir múltiples rutas posibles entre las ciudades y regiones, multiplicadas por las tres alternativas de transporte y los tres años de planificación, la cantidad de variables continuas ($X_{crta}$) crece exponencialmente, obligando al solver a evaluar una inmensa cantidad de combinaciones para satisfacer la demanda.
    
    \item \textbf{Peso de los costos de apertura}: Aunque los costos de apertura se suman una sola vez como constantes asociadas a una variable discreta, influyen drásticamente en la solución. Por ejemplo, al evaluar a Santiago como reemplazo de Antofagasta, Santiago presenta mejores parámetros operativos (menor costo variable y un transporte más barato hacia la región R1, además de un menor costo fijo anual). Sin embargo, la constante de apertura de Santiago penaliza la función objetivo en aproximadamente 85 millones extra en comparación a Antofagasta, provocando que el algoritmo la descarte de la solución óptima.

    \item \textbf{Escalamiento temporal de los costos de producción}: El parámetro de costo variable es crítico debido a su comportamiento escalable en el tiempo. A diferencia de las constantes de apertura, este costo se multiplica directamente por el volumen de demanda, el cual va en aumento año tras año. Si la planificación del modelo fuera mayor a tres años, la función objetivo se vería dominada por este factor, obligando al solver a priorizar la instalación de plantas en los nodos con el menor costo variable posible.
\end{itemize}
\end{document}
